\ifdefined\gkver\else\def\gkver{student}\fi
\documentclass[\gkver]{gaokaozhenti}
\begin{document}

\chapter{2007年上海理综}
%% paperType: 理综物理部分

\begin{enumerate}
\item
%% number: 4
%% paperName: 2007年上海理综第 4 题 3 分
%% typeId: 1
%% type: 单选题
%% chapter: 直线运动
%% point: 匀速直线运动
%% method: 无
%% score: 3
%% degree: 850
%% duplicateId: 0
%% body:
汉语成语短小精悍、言简意赅，蕴涵着许多科学知识。例如“一寸光阴一寸金”中的“寸”出自古代计时器“日晷”。右表成语对应了不同的科学知识，其中正确的是\xzanswer{B}

\begin{center}
\begin{tblr}{|c|c|}
\hline
成语 & 蕴涵的科学知识 \\
\hline
①刻舟求剑 & 沉浮原理 \\
\hline
②海市蜃楼 & 相对运动 \\
\hline
③种瓜得瓜 & 遗传现象 \\
\hline
④炉火纯青 & 光色测温 \\
\hline
\end{tblr}
\end{center}

\fourchoices[answer=B]
{②③}
{③④}
{①④}
{①③}

%% answer: B
%% memo:
\memoanswer{
本题原卷及材料中未提供详解，待补充。
}

\item
%% number: 5
%% paperName: 2007年上海理综第 5 题 3 分
%% typeId: 1
%% type: 单选题
%% chapter: 功和能
%% point: 功能中的图象
%% method: 无
%% score: 3
%% degree: 850
%% duplicateId: 0
%% body:
如图所示显示跳水运动员从离开跳板到入水前的过程。下列正确反映运动员的动能 $E_{k}$ 随时间 $t$ 变化的曲线图是（忽略空气阻力）\xzanswer{C}

\onepicture[w=1.8cm]{figs/05.png}

\fourchoices[answer=C, ispicture=true, h=2.4cm]
{figs/05a.png}
{figs/05b.png}
{figs/05c.png}
{figs/05d.png}

%% answer: C
%% memo:
\memoanswer{
本题原卷及材料中未提供详解，待补充。
}

\item
%% number: 6
%% paperName: 2007年上海理综第 6 题 3 分
%% typeId: 1
%% type: 单选题
%% chapter: 力物体的平衡
%% point: 力矩平衡
%% method: 无
%% score: 3
%% degree: 850
%% duplicateId: 0
%% body:
右图是居民小区内常见的健身器材。使用时手可以握在①或②处，脚可以踩在④或⑤处。手脚一起用力时，器械前半部绕支点③转动。下列最费力的方式是\xzanswer{D}

\onepicture[w=3.2cm]{figs/06.png}

\fourchoices[answer=D]
{手握①，同时脚踩⑤}
{手握①，同时脚踩④}
{手握②，同时脚踩⑤}
{手握⑨，同时脚踩④}

%% answer: D
%% memo:
\memoanswer{
本题原卷及材料中未提供详解，待补充。
}

\item
%% number: 7
%% paperName: 2007年上海理综第 7 题 3 分
%% typeId: 1
%% type: 单选题
%% chapter: 气体性质
%% point: 气体的状态参量
%% method: 无
%% score: 3
%% degree: 650
%% duplicateId: 0
%% body:
有一压力锅，锅盖上的排气孔截面积约为 $7.0\times10^{-6}\Umq$，限压阀村重为 $0.7\UN$。使用该压力锅煮水消毒，根据下列水的沸点与气压的关系的表格，分析可知压力锅内的最高水温约为\xzanswer{C}（大气压强为 $1.01\times10^{5}\UPa$）

\begin{center}
\small
\begin{tblr}{|c|c|c|c|c|c|c|c|c|c|c|}
\hline
$p$（$\times10^{5}\UPa$） & 1.01 & 1.43 & 1.54 & 1.63 & 1.73 & 1.82 & 1.91 & 2.01 & 2.12 & 2.21 \\
\hline
$t$（$\UdC$） & 100 & 110 & 112 & 114 & 116 & 118 & 120 & 122 & 124 & 126 \\
\hline
\end{tblr}
\end{center}

\onepicture[w=2.6cm, align=right]{figs/07.png}

\fourchoices[answer=C]
{$100\UdC$}
{$112\UdC$}
{$122\UdC$}
{$124\UdC$}

%% answer: C
%% memo:
\memoanswer{
本题原卷及材料中未提供详解，待补充。
}

\item
%% number: 8
%% paperName: 2007年上海理综第 8 题 3 分
%% typeId: 1
%% type: 单选题
%% chapter: 力物体的平衡
%% point: 受力分析
%% method: 无
%% score: 3
%% degree: 750
%% duplicateId: 0
%% body:
在热气球下方开口处燃烧液化气，使热气球内部气体温度升高，热气球开始离地，徐徐升空。分析这一过程，下列表述正确的是\xzanswer{B}

①气球内的气体密度变小，所受重力也变小

②气球内的气体密度不变，所受重力也不变

③气球所受浮力变大

④气球所受浮力不变

\onepicture[w=2.2cm]{figs/08.png}

\fourchoices[answer=B]
{①③}
{①④}
{②③}
{②④}

%% answer: B
%% memo:
\memoanswer{
本题原卷及材料中未提供详解，待补充。
}

\item
%% number: 9
%% paperName: 2007年上海理综第 9 题 3 分
%% typeId: 1
%% type: 单选题
%% chapter: 曲线运动
%% point: 天体运动
%% method: 无
%% score: 3
%% degree: 650
%% duplicateId: 0
%% body:
太阳系八大行星公转轨道可近似看作圆轨道，“行星公转周期的平方”与“行星与太阳的平均距离的三次方”成正比。地球与太阳之间平均距离约为 1.5 亿千米，结合下表可知，火星与太阳之间的平均距离约为\xzanswer{B}

\begin{center}
\begin{tblr}{|c|c|c|c|c|c|c|}
\hline
 & 水星 & 金星 & 地球 & 火星 & 木星 & 土星 \\
\hline
公转周期（年） & 0.241 & 0.615 & 1.0 & 1.88 & 11.86 & 29.5 \\
\hline
\end{tblr}
\end{center}

\fourchoices[answer=B]
{1.2 亿千米}
{2.3 亿千米}
{4.6 亿千米}
{6.9 亿千米}

%% answer: B
%% memo:
\memoanswer{
本题原卷及材料中未提供详解，待补充。
}

\item
%% number: 11
%% paperName: 2007年上海理综第 11 题 3 分
%% typeId: 1
%% type: 单选题
%% chapter: 原子和原子核
%% point: 半衰期
%% method: 无
%% score: 3
%% degree: 750
%% duplicateId: 0
%% body:
放射性同位素 \ce{^{14}C} 可用来推算文物的“年龄”。\ce{^{14}C} 的含量每减少一半要经过约 5730 年。某考古小组挖掘到一块动物骨骼，经测定 \ce{^{14}C} 还剩余 $1/8$，推测该动物生存年代距今约为\xzanswer{A}

\fourchoices[answer=A]
{$5730\times3$ 年}
{$5730\times4$ 年}
{$5730\times6$ 年}
{$5730\times8$ 年}

%% answer: A
%% memo:
\memoanswer{
本题原卷及材料中未提供详解，待补充。
}

\item
%% number: 42
%% paperName: 2007年上海理综第 42 题 2 分
%% typeId: 3
%% type: 填空题
%% chapter: 曲线运动
%% point: 人造地球卫星
%% method: 无
%% score: 2
%% degree: 750
%% duplicateId: 0
%% body:
据报道，中国第一颗人造“探月”卫星将于 2007 年下半年发射。有人建议，从节省火箭燃料考虑，将发射基地从西昌移至海南。理由是\tkanswer[4cm]{海南的自转线速度大}。

%% answer: 
\jdanswer{
海南的自转线速度大
}
%% memo:
\memoanswer{
本题原卷及材料中未提供详解，待补充。
}

\item
%% number: 43
%% paperName: 2007年上海理综第 43 题 3 分
%% typeId: 1
%% type: 单选题
%% chapter: 曲线运动
%% point: 万有引力定律
%% method: 无
%% score: 3
%% degree: 750
%% duplicateId: 0
%% body:
据悉，上海为配合“嫦娥一号”工程，启动了“吴刚计划”，正式开始研制“登月车”。根据月球与地球的环境差异，研制“登月车”时应考虑的月球表面因素是\xzanswer{C}

①昼夜温差大　　②宇宙射线较弱　　③高低起伏大　　④重力加速度小

\fourchoices[answer=C]
{①②③}
{①②④}
{①③④}
{②③④}

%% answer: C
%% memo:
\memoanswer{
本题原卷及材料中未提供详解，待补充。
}

\item
%% number: 44
%% paperName: 2007年上海理综第 44 题 2 分
%% typeId: 3
%% type: 填空题
%% chapter: 气体性质
%% point: 能源的开发与利用
%% method: 无
%% score: 2
%% degree: 850
%% duplicateId: 0
%% body:
利用新能源发电的形式很多（见下图），其中 A 表示潮汐能发电，B 表示地热能发电，C 表示\tkanswer[1cm]{风}能发电，D 表示\tkanswer[1.2cm]{太阳}能发电。

\fourpicture[numA=A, numB=B, numC=C, numD=D, labelprefix={}, w=2.8cm]
{figs/44a.png}
{figs/44b.png}
{figs/44c.png}
{figs/44d.png}

%% answer: 
\jdanswer{
风；太阳
}
%% memo:
\memoanswer{
本题原卷及材料中未提供详解，待补充。
}

\item
%% number: 45
%% paperName: 2007年上海理综第 45 题 9 分
%% typeId: 3
%% type: 填空题
%% chapter: 气体性质
%% point: 能源的开发与利用
%% method: 无
%% score: 9
%% degree: 650
%% duplicateId: 0
%% body:
我国自古有“昼涨称潮，夜涨称汐”的说法。潮汐主要是由太阳和月球对海水的引力造成的，以月球对海水的引力为主。

\begin{enumerate}
	\item
	世界两大观潮胜地，一处是亚马孙河北河口，另一处是我国的\tkanswer[2.5cm]{钱塘江}入海口。

	\item
	右图\ref{2007上海理综45a}、\ref{2007上海理综45b}是某类潮汐发电示意图。涨潮时开闸，水由通道进入海湾水库蓄水，待水面升至最高点时关闭闸门（见图\ref{2007上海理综45a}）。当落潮时，开闸放水发电（见图\ref{2007上海理综45b}）。设海湾水库面积为 $5.0\times10^{8}\Umq$，平均潮差为 $3.0\Um$，一天涨落潮两次，发电的平均能量转化率为 $10\%$，发电的平均功率约为\xzanswer{B}。（$\rho_{\text{海水}}$ 取 $1.0\times10^{3}\Ukgmc$，$g$ 取 $10\Umsq$）

	\fourchoices[answer=B]
	{$2.6\times10^{4}\UkW$}
	{$5.2\times10^{4}\UkW$}
	{$2.6\times10^{5}\UkW$}
	{$5.2\times10^{5}\UkW$}

	\twopicture[numA=a, numB=b, labelprefix={图}, labelA=2007上海理综45a, labelB=2007上海理综45b, align=right]
	{figs/45a.png}
	{figs/45b.png}

	\item
	右图为双水库潮汐电站原理示意图。两个水库之间始终保持着水位差，可以全天发电。涨潮时，闸门的开关情况是\tkanswer[2cm]{A关B开}；落潮时，闸门的开关情况是\tkanswer[2cm]{B关A开}。从能量的角度说，该电站是将海水的\tkanswer[1.2cm]{势能}转化为水轮机的动能，再推动发电机发电。

	\onepicture[align=right, w=4.5cm]{figs/45c.png}
\end{enumerate}

%% answer: 
\jdanswer{
\begin{enumerate}
	\item
	钱塘江

	\item
	B

	\item
	A关B开；B关A开；势能

\end{enumerate}
}
%% memo:
\memoanswer{
本题原卷及材料中未提供详解，待补充。
}

\item
%% number: 46
%% paperName: 2007年上海理综第 46 题 10 分
%% typeId: 3
%% type: 填空题
%% chapter: 原子和原子核
%% point: 核裂变
%% method: 无
%% score: 10
%% degree: 650
%% duplicateId: 0
%% body:
潮汐能属于无污染能源，但能量的转化率较低，相比之下，核能是一种高效的能源。

\begin{enumerate}
	\item
	在核电站中，为了防止放射性物质泄露，核反应堆有三道防护屏障：燃料包壳，压力壳和安全壳（见图\ref{2007上海理综46a}）。结合图\ref{2007上海理综46b}可知，安全壳应当选用的材料是\tkanswer[2cm]{混凝土}。

	\item
	核反应堆中的核废料具有很强的放射性，目前常用的处理方法是将其装入特制的容器中，然后\xzanswer{D}

	\fourchoices[answer=D]
	{沉入海底}
	{放至沙漠}
	{运到月球}
	{深埋地下}

	\threepicture[numstyle=arabic, labelprefix={图}, labelA=2007上海理综46a, labelB=2007上海理综46b, labelC=2007上海理综46c, align=right, w=3.2cm]
	{figs/46a.png}
	{figs/46b.png}
	{figs/46c.png}

	\item
	图\ref{2007上海理综46c}是用来监测工作人员受到辐射情况的胸章，通过照相底片被射线感光的区域，可以判断工作人员受到何种辐射。当胸章上 $1\Umm$ 铝片和 $3\Umm$ 铝片下的照相底片被感光，而铅片下的照相底片未被感光时，结合图\ref{2007上海理综46b}分析可知工作人员受到了\tkanswer[1cm]{$\beta$}射线的辐射；当所有照相底片被感光时，工作人员受到了\tkanswer[1cm]{$\gamma$}射线的辐射。
\end{enumerate}

%% answer: 
\jdanswer{
\begin{enumerate}
	\item
	混凝土

	\item
	D

	\item
	$\beta$；$\gamma$

\end{enumerate}
}
%% memo:
\memoanswer{
本题原卷及材料中未提供详解，待补充。
}

\item
%% number: 47
%% paperName: 2007年上海理综第 47 题 3 分
%% typeId: 1
%% type: 单选题
%% chapter: 气体性质
%% point: 能源的开发与利用
%% method: 无
%% score: 3
%% degree: 850
%% duplicateId: 0
%% body:
氢能是一种既高效又干净的新能源，发展前景良好，用氢作能源的燃料电池汽车倍受青睐。我国拥有完全自主知识产权的氢燃料电池轿车“超越三号”，已达到世界先进水平，并加快向产业化的目标迈进。氢能具有的优点包括\xzanswer{A}

①原料来源广　　②易燃烧、热值高　　③储存方便　　④制备工艺廉价易行

\fourchoices[answer=A]
{①②}
{①③}
{③④}
{②④}

%% answer: A
%% memo:
\memoanswer{
本题原卷及材料中未提供详解，待补充。
}

\end{enumerate}

\end{document}
